%
% 7-fluegel.tex -- Einfluss der Flügelgeometrie auf die Eigenwerte
%
% (c) 2026 Stephanie Märklin, OST Ostschweizer Fachhochschule
%
% !TEX root = ../../buch.tex
% !TEX encoding = UTF-8
%
\section{Einfluss der Flügelgeometrie auf die Eigenwerte
\label{aircraft:section:fluegel}}
\kopfrechts{Flügelgeometrie}
Die Systemmatrizen des Propellerflugzeugs und des Business Jets in
Abschnitt~\ref{aircraft:section:6-beispiel} haben dieselbe Form.
Die beiden Flugzeuge unterscheiden sich nur in den Werten der Derivative.
Diese Werte entstehen aus der Geometrie des Flugzeugs.
Wir betrachten die drei Derivative $L_\beta$, $N_\beta$ und $N_r$.
Diese drei bestimmen die Spiralbewegung und die Dutch Roll, und der
Konstrukteur legt sie über die Form der Flügel und des Seitenleitwerks fest.

Die \emph{V-Form} (dihedral) ist der Winkel, um den die Flügel gegenüber der
Querachse nach oben geneigt sind.
Schiebt das Flugzeug, so trifft die Luft den Flügel auf der Seite, zu der das
Flugzeug schiebt, unter einem grösseren Anstellwinkel als den anderen Flügel.
Dieser Flügel erzeugt mehr Auftrieb, und das Flugzeug rollt von dieser Seite
weg.
Dieses Rollmoment infolge Schiebens ist das Derivativ $L_\beta$ aus
Abschnitt~\ref{aircraft:section:system}.
Mit wachsender V-Form wächst der Betrag von $L_\beta$.

Die \emph{Pfeilung} (sweep) ist der Winkel, um den die Flügel gegenüber der
Querachse nach hinten geneigt sind.
Beim Schieben steht der Flügel auf der Seite, zu der das Flugzeug schiebt,
weniger schräg zur Anströmung als der andere Flügel und erzeugt ebenfalls mehr
Auftrieb.
Die Pfeilung vergrössert den Betrag von $L_\beta$ ebenso wie die V-Form.

Das \emph{Seitenleitwerk} (vertical tail) bestimmt die beiden
Derivative $N_\beta$ und $N_r$.
Schiebt das Flugzeug, so wird das Seitenleitwerk schräg angeströmt und erzeugt
ein Giermoment, das die Nase in die Anströmung dreht.
Dieses Giermoment infolge Schiebens ist die Richtungsstabilität $N_\beta$.
Giert das Flugzeug, so bewegt sich das Seitenleitwerk quer zur Flugrichtung und
erzeugt ein Giermoment gegen die Drehung.
Dieses Giermoment infolge Gierens ist die Gierdämpfung $N_r$.
Ein grösseres Seitenleitwerk vergrössert die Beträge von $N_\beta$ und $N_r$
zugleich.
Die Dutch Roll wird begünstigt, wenn der Betrag von $L_\beta$ gross ist im
Vergleich zu $N_\beta$.

Wie die Derivative die Eigenwerte festlegen, zeigt die Spiralbewegung am
einfachsten, denn ihr Eigenwert ist reell und lässt sich näherungsweise in
geschlossener Form angeben.
Nelson \cite{aircraft:nelson} gibt die Näherung
\begin{equation}
    \lambda_{\text{S}}
    \approx
    \frac{L_\beta N_r - L_r N_\beta}{L_\beta}
    \label{aircraft:eqn:spiralnaeherung}
\end{equation}
an, wobei $L_r$ das Rollmoment infolge Gierens ist.
Bei üblichen Vorzeichen sind $L_\beta$ und $N_r$ negativ, $L_r$ und $N_\beta$
positiv.
Der Nenner ist also negativ, und der Eigenwert der Spiralbewegung ist genau dann
negativ, wenn
\begin{equation}
    L_\beta N_r - L_r N_\beta > 0
    \label{aircraft:eqn:spiralbedingung}
\end{equation}
gilt.
Tatsächlich ist für das Propellerflugzeug mit den Werten aus Nelson Tabelle~5.2
\[
    L_\beta N_r - L_r N_\beta
    =
    (-16{,}02)\cdot(-0{,}76) - 2{,}19\cdot 4{,}49
    =
    12{,}18 - 9{,}83
    >
    0.
\]
Der Eigenwert der Spiralbewegung ist nach
Abschnitt~\ref{aircraft:section:6-beispiel} negativ, wie es die Bedingung
verlangt.
Die Bedingung \eqref{aircraft:eqn:spiralbedingung} wird umso leichter erfüllt,
je grösser die Beträge von $L_\beta$ und $N_r$ sind.
Ob die Spiralbewegung abklingt oder anwächst, entscheidet das Verhältnis der
vier Derivative und damit die Konstruktion.

Nelsons Näherung für die Dutch Roll vernachlässigt die Rollbewegung und enthält
$L_\beta$ deshalb nicht.
Nelson \cite[Abb.~5.14]{aircraft:nelson} zeigt stattdessen, wie die Eigenwerte
in der komplexen Ebene wandern, wenn der Betrag von $L_\beta$ wächst.
Der Eigenwert der Spiralbewegung wandert nach links, das Eigenwertpaar der
Dutch Roll wandert nach rechts.
Eine grosse V-Form stabilisiert die Spiralbewegung und schwächt die Dämpfung der
Dutch Roll.
Eine grosse Richtungsstabilität $N_\beta$ wirkt umgekehrt, denn sie erhöht die
Frequenz der Dutch Roll und zieht den Eigenwert der Spiralbewegung zum Ursprung.
Der Konstrukteur muss beide Eigenwerte zugleich in der linken Halbebene halten.
Die Gierdämpfung $N_r$ vergrössert die Dämpfung beider Bewegungen, lässt sich
durch ein grösseres Seitenleitwerk aber nur zusammen mit $N_\beta$ vergrössern.

Der Vergleich in Abschnitt~\ref{aircraft:section:6-beispiel} ist mit diesem
Zusammenhang verträglich.
Der Business Jet hat gepfeilte Flügel, das Propellerflugzeug gerade.
Die Dutch Roll des Business Jets ist mit $\sigma = -0{,}116$ schwächer gedämpft
als die des Propellerflugzeugs mit $\sigma = -0{,}487$.

Die Geometrie legt die Derivative fest, und die Derivative legen die
Eigenwerte fest.
Ob ein Flugzeug stabil ist, wird damit beim Entwurf entschieden.